\section{Why defaults recur}
\label{sec:hypotheses}

Surface patterns alone do not reveal their causes. The explanations below are
hypotheses about pressures that may favor portable language and local
regularity. They are useful because they generate different predictions, not
because the prose alone proves any one of them.

\subsection{Local prediction rewards local plausibility}

During pretraining, a language model learns to predict each next token from the
preceding context. Later training changes its behavior, but familiar phrases
and locally smooth continuations may remain attractive even when the paragraph
requires a less common structure. A stock transition can therefore fit its
immediate sentence while weakening the larger argument.

This account helps explain two apparently opposite defaults. Repeated formulas
preserve familiar shapes, while synonym cycling avoids lexical repetition by
changing the name of a stable object. In both cases, a locally plausible choice
forces the reader to recover the paragraph-level relation.

\subsection{Missing reader context leaves priorities unclear}

A prompt can name a task without identifying the person who will use the
answer. In that setting, the model has little basis for deciding which
background to omit or which consequence matters. It may not know whether an
objection deserves space. Generic helpfulness then fills the gap with broad
explanation and visible signposting; the answer tends toward balanced coverage
and portable management language.

The reader hypothesis predicts that specificity should improve the selection
of detail, not merely the tone. A response written for a reviewer with a
pending decision ought to emphasize different evidence from one written for a
student encountering the mechanism for the first time.

\subsection{Training may reward visible completeness}

Instruction tuning and preference optimization may favor features that
evaluators readily recognize as helpful. Headings, long lists, summaries,
caveats, and offers to continue make a chat response look thorough. Transferred
into a finished article, the same habits can interrupt continuity or append
material after the argument has ended.

Mixed training data may also favor a middle register acceptable across many
genres. Code documentation and issue trackers contribute terse inventory
prose, while business writing supplies a managerial register. Essays and
reference works add other conventions. Without a strong cue about voice and
genre, their most portable features can combine into prose that is competent
but weakly situated.

The three pressures can reinforce one another. Local prediction supplies the
familiar phrase, post-training rewards its recognizable helpfulness, and
missing reader context gives the model no reason to remove it. The resulting
passage can satisfy the conventions of a helpful answer without resolving the
problem that brought the reader to it.
